\title{Assignment 5: CS 215}
\author{}
\date{Due: 7th November before 11:55 pm}

\documentclass[11pt]{article}

\usepackage{amsmath,xcolor}
\usepackage{amssymb}
\usepackage{hyperref}
\usepackage{soul}
\usepackage[margin=0.4in]{geometry}
\begin{document}
\maketitle

\textbf{Remember the honor code while submitting this (and every other) assignment. You may discuss broad ideas with other students or ask me for any difficulties, but the code you implement and the answers you write must be your own. We will adopt a \textbf{zero-tolerance policy} against any violation.}
\\
\\
\textbf{Submission instructions:} Follow the instructions for the submission format and the naming convention of your files from the submission guidelines file in the homework folder. However, please do \emph{not} submit the face image databases in your zip file that you will upload on moodle. Please see \textsf{assignment4\_CS215.rar}. Upload the file on moodle \emph{before} 11:55 pm on 7th November. \textbf{\textcolor{red}{Since this is the last day of classes, the cutoff date/time this time is also 11:55 pm on 7th November. No late assignments will be accepted after this time.}} Please preserve a copy of all your work until the end of the semester.  

\begin{enumerate}
\item The aim of this exercise is to help you understand the mathematics behind PCA more deeply. Do as directed: \textsf{[5+5+5+5=20 points]}
\begin{enumerate}
\item Prove that the covariance matrix in PCA is symmetric and positive semi-definite.
\item Prove that the eigenvectors of a symmetric matrix are orthonormal.
\item Consider a dataset of some $N$ vectors in $d$ dimensions given by $\{\boldsymbol{x_i}\}_{i=1}^N$ with mean vector $\boldsymbol{\bar{x}}$. Note that each $\boldsymbol{x_i} \in \mathbb{R}^d$ and also $\boldsymbol{\bar{x}} \in \mathbb{R}^d$.
Suppose that only $k$ eigenvalues of the corresponding covariance matrix are large and the remaining are very small in value. Let $\boldsymbol{\tilde{x}_i}$ be an approximation to $\boldsymbol{x_i}$ of the form $\boldsymbol{\tilde{x}_i} = \boldsymbol{\bar{x}} + \sum_{l=1}^k \boldsymbol{V_l} \alpha_{il}$ where $\boldsymbol{V_l}$ stands for the $l$th eigenvector (it has $d$ elements) and $\alpha_{il}$ (it is a scalar) stands for the $l$th eigencoefficient of $\boldsymbol{x_i}$. Argue why the error $\frac{1}{N} \sum_{i=1}^N \|\boldsymbol{\tilde{x}_i} - \boldsymbol{x_i}\|^2_2$ will be small. What will be the value of this error in terms of the eigenvalues of the covariance matrix?   
\item Consider two uncorrelated zero-mean jointly Gaussian random variables $(X_1, X_2)$. Let $X_1$ belong to a Gaussian distribution with variance 100 and $X_2$ belong to a Gaussian distribution with variance 1. What are the principal components of $(X_1, X_2)$? If the variance of $X_1$ and $X_2$ were equal, what are the principal components? 
\end{enumerate}

\item In this part, you will implement a mini face recognition system. For this, read the instructions below and do as directed:
\begin{enumerate}
\item Download the ORL face database from the homework folder. It contains 40 sub-folders, one for each of the 40 subjects/persons. For each person, there are ten images in the appropriate folder named 1.pgm to 10.pgm. The images are of size 92 by 110 each. 

\item Each image is in the pgm format which you can read using \texttt{imread} in MATLAB. After reading, please convert the images into double using \begin{verbatim}im = imread('A.pgm'); im = double(im);\end{verbatim}. You can view/read the images in this format, either through MATLAB or through image viewers like IrfanView on Windows, or xv/display/gimp on Unix. Though the face images are in different poses, expressions and facial accessories, they are all roughly aligned (the eyes are in roughly similar locations in all images). 

\item (*) For this first part of the assignment, you will work with the images of the first 32 people (numbers from 1 to 32). For each person, you will include the first six images in the training set (that is the first 6 images that appear in a directory listing as produced by the \textsf{dir} function of MATLAB) and the remaining four images in the testing set. 

\item You should implement the recognition system by using the \textsf{eig} or \textsf{eigs} function of MATLAB on the matrix $\boldsymbol{L}$ which has size $n \times n$ where $n$ is the total number of training images. Do \textbf{not} use the covariance matrix as it will be too large in size as explained in class, and as seen in a problem in HW4. Record the recognition rate using squared difference between the eigen-coefficients while testing on all the images in the test set, for $k \in \{1,2,3,5,10,15,20,30,50,75,100,150,170\}$. Plot the rates in your report in the form of a graph.

\item Repeat the same experiment (using just the \textsf{eig} or \textsf{eigs} routine) on the Yale Face database from the homework folder. This database contains about 64 images each of 38 individuals (\textit{labeled from 1 to 39, with number 14 missing; some folders have slightly less than 64 images}). Each image is in pgm format and has size 192 by 168. The images are taken under different lighting conditions but in the same pose. 

\item Take the first 40 images of every person for training and test on the remaining 24 images (that is the first 40 images that appear in a directory listing as produced by the \textsf{dir} function of MATLAB). 

\item For both ORL and Yale databases, plot in your report the recognition rates for $k \in \{1,2,3,5,10,15,20,30,50,60, 65,75,100,200,300,500,1000\}$ based on (a) the squared difference between all the eigencoefficients and (b) the squared difference between all \emph{except} the three eigencoefficients corresponding to the eigenvectors with the three largest eigenvalues. 

\item Display in your report the reconstruction of any one face image from the ORL database using $k \in \{2,10,20,50,75,100,125, 150,175\}$ values. Plot the 25 eigenvectors (eigenfaces) corresponding to the 25 largest eigenvalues using the subplot or subimage commands in MATLAB. \textsf{[40 points]}
\end{enumerate}

\end{enumerate}



\end{document}